\section{Cloud-native and traditional engines on a single node}
\label{discussion:single-node-cloud-native-vs-traditional}

Section~\ref{results:single-machine-query-performance} reports the single-machine orderings on wall-clock time, network transfer, and operational cost but does not attribute them to a cause. Reading those orderings against the mechanisms of Chapter~\ref{ch:background} explains why the cloud-native combination of GeoParquet and DuckDB behaved as it did against the traditional PostGIS and Shapefile alternatives, the comparison \acrshort{rq}1 poses. The three dimensions are treated separately for two reasons: the engines that lead on one do not lead on the others, and the host-boundary byte count carries a different meaning for each configuration, which resists a single ranking.

On run-time the single-machine ordering tracks the spatial-indexing capability each stack natively exposes (Figure~\ref{fig:single-machine-wall-clock-time}, Table~\ref{tab:results-achieved-n}). The order-of-magnitude gap between PostGIS and the GeoParquet path is compounded by the decoupled read model rather than originating in it. PostGIS answers a selective query from a warm, server-resident process whose geometries are backed by a persistent spatial index (Section~\ref{background:geometry-and-indexing:spatial-indexing}), so it touches a small, resident working set. DuckDB over GeoParquet instead reads its input across a decoupled storage boundary, and that separation imposes a per-request latency and a processing cost on every scan that a co-located engine does not pay \parencite{dageville2016_snowflake_elastic_data_warehouse,pang2024_storage_disaggregated_databases} (Section~\ref{background:cloud-native-data-architecture:decoupled-storage-and-compute}). At the sub-second scale of these queries the round-trips to object storage dominate the interval \parencite{fielding2022_rfc9110_http_semantics} (Section~\ref{background:cloud-native-data-architecture:object-storage-and-http-access}), so the cloud-native read model did not convert into a wall-clock win against a warm, indexed local engine. Because each engine ran in its native configuration, this ordering reflects indexing capability and storage locality together with the engine, and so reads as a property of the stacks as deployed rather than of the engines in isolation, a bound taken up in Section~\ref{discussion:reach-and-limits}. How DuckDB reaches its input is a property of that read model, not a transfer advantage, as the byte signatures that follow show.

% Academic rewrite
Network transfer cannot be read as a single ranking, because the host-boundary byte count denotes a different quantity for each configuration: the input scanned from object storage for DuckDB, the returned result set for PostGIS, and a near-zero local read for the Shapefile path (Section~\ref{results:consistency-of-winners}). Read per architecture, the received volumes follow from how each engine reaches its data (Figure~\ref{fig:single-machine-bytes-transferred}). DuckDB received a near-constant payload across the three query patterns, so its per-row transfer fell as the result cardinality grew (Figure~\ref{fig:bytes-vs-cardinality}). That is the signature of range-request access with column and row-group pruning. The GeoParquet layout produces it: rows sit in independently readable row groups, columnar within each group, behind a footer index of per-row-group statistics, so a query reads a footprint set by the columns and row groups it touches rather than by the number of rows it returns \parencite{ogc2024_geoparquet_specification_v110,lamb2022_querying_parquet_millisecond_latency,maso2023_ogc_cog_standard,fielding2022_rfc9110_http_semantics,duckdb2025_parquet_overview} (Figure~\ref{fig:geoparquet-file-format}, Section~\ref{background:cloud-native-data-architecture:object-storage-and-http-access}, Section~\ref{background:spatial-data-formats-and-storage-models:file-formats:geoparquet}). PostGIS, by contrast, returns the result set itself, so its received volume rose with the row count. The two read models account for the divergence, which is no verdict on either engine.

The clearest evidence that transfer is not one axis is that the larger mover reverses by pattern. On the point-in-polygon lookup at the small tier DuckDB receives the larger payload, while on the bounding-box filter the order inverts and PostGIS receives more than an order of magnitude more, because the filter returns many rows while DuckDB's footprint stays near-constant (Figure~\ref{fig:single-machine-bytes-transferred}, Figure~\ref{fig:bytes-vs-cardinality}). A quantity whose ordering depends on the query pattern cannot place the configurations on one common byte scale. The direction of the traffic carries a further signal. For both engines received bytes exceeded sent bytes by orders of magnitude (Figure~\ref{fig:single-machine-bytes-directional}), and that split records where the work is performed (Section~\ref{research-design-and-methodology:metrics-and-cost-model:primary-metrics}). For PostGIS the received bytes are the result set, so the server did the spatial work and returned an answer. For DuckDB the received bytes are the scanned input, so the client pulled data in order to compute over it. The asymmetry locates computation; it does not score one configuration above another on transfer.

Cost and time did not move together. DuckDB recorded the lower operational cost on the small-tier k-nearest-neighbor search and bounding-box filter even though it lost to PostGIS on run-time in those cells (Figure~\ref{fig:single-machine-operational-cost}, Figure~\ref{fig:single-machine-cost-time-quadrant}). Under per-resource metering, where each resource is billed on its own dimension \parencite{mell2011_nist_definition_of_cloud_computing} (Section~\ref{background:cloud-native-data-architecture:cloud-cost-model}), the charge follows the resource a workload actually consumes rather than its elapsed time. At these data sizes the compute term carried the charge for both single-node engines. The byte signatures of the preceding paragraphs did not enter this ranking. Transfer reaches operational cost only through the egress term of the cost model (Equation~\ref{eq:cost-model}, Section~\ref{research-design-and-methodology:metrics-and-cost-model:cost-model}), and that term registered zero for every single-node configuration (Section~\ref{results:single-machine-query-performance:operational-cost}). The large differences in bytes received therefore produced no cost difference at the single-node scale tested. How far this holds is bounded by that scale: a configuration billed for cross-region egress, or one scanning enough input to move the compute term, would couple transfer to cost in a way these measurements do not reach.

% Academic rewrite
The remaining single-node result is the GeoPandas over Shapefile path. A \texttt{.qix} quadtree index was built over the bundle during setup (Section~\ref{system-architecture-and-implementation:data-layer:shapefile-bundle}), so the path is not unindexed; even so, the k-nearest-neighbor search is the pattern on which it falls furthest behind, the slowest configuration measured there by more than an order of magnitude against DuckDB (Figure~\ref{fig:single-machine-cross-pattern}). A quadtree, like any containment index, prunes range and intersection predicates but does not order by distance without a bounding-box predicate \parencite{guttman1984_rtree_dynamic_index} (Section~\ref{background:geometry-and-indexing:spatial-indexing}), so the nearest-neighbor search gains nothing from it, just as DuckDB's row-group pruning does not serve that pattern. The gap to DuckDB there is consistent with the absence of a query engine on the GeoPandas path, where the search resolves through an in-memory pass over deserialized geometries rather than a vectorized scan-and-sort. On the point-in-polygon and bounding-box patterns the same path stayed competitive. This result belongs to the GeoPandas and Shapefile combination as exercised, not to the Shapefile format in isolation, because the format was reached only through that library. That distinction qualifies how far the finding generalizes, a point taken up in Section~\ref{discussion:reach-and-limits}.